% Options for packages loaded elsewhere
\PassOptionsToPackage{unicode}{hyperref}
\PassOptionsToPackage{hyphens}{url}
\PassOptionsToPackage{dvipsnames,svgnames,x11names}{xcolor}
\documentclass[
  11pt,
]{article}
\usepackage{xcolor}
\usepackage[margin=2cm]{geometry}
\usepackage{amsmath,amssymb}
\setcounter{secnumdepth}{-\maxdimen} % remove section numbering
\usepackage{iftex}
\ifPDFTeX
  \usepackage[T1]{fontenc}
  \usepackage[utf8]{inputenc}
  \usepackage{textcomp} % provide euro and other symbols
\else % if luatex or xetex
  \usepackage{unicode-math} % this also loads fontspec
  \defaultfontfeatures{Scale=MatchLowercase}
  \defaultfontfeatures[\rmfamily]{Ligatures=TeX,Scale=1}
\fi
\usepackage{lmodern}
\ifPDFTeX\else
  % xetex/luatex font selection
    \setmonofont[Scale=0.8,BoldFont=DejaVuSansMono-Bold.ttf]{DejaVuSansMono.ttf}
\fi
% Use upquote if available, for straight quotes in verbatim environments
\IfFileExists{upquote.sty}{\usepackage{upquote}}{}
\IfFileExists{microtype.sty}{% use microtype if available
  \usepackage[]{microtype}
  \UseMicrotypeSet[protrusion]{basicmath} % disable protrusion for tt fonts
}{}
\makeatletter
\@ifundefined{KOMAClassName}{% if non-KOMA class
  \IfFileExists{parskip.sty}{%
    \usepackage{parskip}
  }{% else
    \setlength{\parindent}{0pt}
    \setlength{\parskip}{6pt plus 2pt minus 1pt}}
}{% if KOMA class
  \KOMAoptions{parskip=half}}
\makeatother
\usepackage{color}
\usepackage{fancyvrb}
\newcommand{\VerbBar}{|}
\newcommand{\VERB}{\Verb[commandchars=\\\{\}]}
\DefineVerbatimEnvironment{Highlighting}{Verbatim}{commandchars=\\\{\}}
% Add ',fontsize=\small' for more characters per line
\newenvironment{Shaded}{}{}
\newcommand{\AlertTok}[1]{\textcolor[rgb]{1.00,0.00,0.00}{\textbf{#1}}}
\newcommand{\AnnotationTok}[1]{\textcolor[rgb]{0.38,0.63,0.69}{\textbf{\textit{#1}}}}
\newcommand{\AttributeTok}[1]{\textcolor[rgb]{0.49,0.56,0.16}{#1}}
\newcommand{\BaseNTok}[1]{\textcolor[rgb]{0.25,0.63,0.44}{#1}}
\newcommand{\BuiltInTok}[1]{\textcolor[rgb]{0.00,0.50,0.00}{#1}}
\newcommand{\CharTok}[1]{\textcolor[rgb]{0.25,0.44,0.63}{#1}}
\newcommand{\CommentTok}[1]{\textcolor[rgb]{0.38,0.63,0.69}{\textit{#1}}}
\newcommand{\CommentVarTok}[1]{\textcolor[rgb]{0.38,0.63,0.69}{\textbf{\textit{#1}}}}
\newcommand{\ConstantTok}[1]{\textcolor[rgb]{0.53,0.00,0.00}{#1}}
\newcommand{\ControlFlowTok}[1]{\textcolor[rgb]{0.00,0.44,0.13}{\textbf{#1}}}
\newcommand{\DataTypeTok}[1]{\textcolor[rgb]{0.56,0.13,0.00}{#1}}
\newcommand{\DecValTok}[1]{\textcolor[rgb]{0.25,0.63,0.44}{#1}}
\newcommand{\DocumentationTok}[1]{\textcolor[rgb]{0.73,0.13,0.13}{\textit{#1}}}
\newcommand{\ErrorTok}[1]{\textcolor[rgb]{1.00,0.00,0.00}{\textbf{#1}}}
\newcommand{\ExtensionTok}[1]{#1}
\newcommand{\FloatTok}[1]{\textcolor[rgb]{0.25,0.63,0.44}{#1}}
\newcommand{\FunctionTok}[1]{\textcolor[rgb]{0.02,0.16,0.49}{#1}}
\newcommand{\ImportTok}[1]{\textcolor[rgb]{0.00,0.50,0.00}{\textbf{#1}}}
\newcommand{\InformationTok}[1]{\textcolor[rgb]{0.38,0.63,0.69}{\textbf{\textit{#1}}}}
\newcommand{\KeywordTok}[1]{\textcolor[rgb]{0.00,0.44,0.13}{\textbf{#1}}}
\newcommand{\NormalTok}[1]{#1}
\newcommand{\OperatorTok}[1]{\textcolor[rgb]{0.40,0.40,0.40}{#1}}
\newcommand{\OtherTok}[1]{\textcolor[rgb]{0.00,0.44,0.13}{#1}}
\newcommand{\PreprocessorTok}[1]{\textcolor[rgb]{0.74,0.48,0.00}{#1}}
\newcommand{\RegionMarkerTok}[1]{#1}
\newcommand{\SpecialCharTok}[1]{\textcolor[rgb]{0.25,0.44,0.63}{#1}}
\newcommand{\SpecialStringTok}[1]{\textcolor[rgb]{0.73,0.40,0.53}{#1}}
\newcommand{\StringTok}[1]{\textcolor[rgb]{0.25,0.44,0.63}{#1}}
\newcommand{\VariableTok}[1]{\textcolor[rgb]{0.10,0.09,0.49}{#1}}
\newcommand{\VerbatimStringTok}[1]{\textcolor[rgb]{0.25,0.44,0.63}{#1}}
\newcommand{\WarningTok}[1]{\textcolor[rgb]{0.38,0.63,0.69}{\textbf{\textit{#1}}}}
\usepackage{longtable,booktabs,array}
\usepackage{calc} % for calculating minipage widths
% Correct order of tables after \paragraph or \subparagraph
\usepackage{etoolbox}
\makeatletter
\patchcmd\longtable{\par}{\if@noskipsec\mbox{}\fi\par}{}{}
\makeatother
% Allow footnotes in longtable head/foot
\IfFileExists{footnotehyper.sty}{\usepackage{footnotehyper}}{\usepackage{footnote}}
\makesavenoteenv{longtable}
\ifLuaTeX
\usepackage[bidi=basic]{babel}
\else
\usepackage[bidi=default]{babel}
\fi
\babelprovide[main,import]{french}
% get rid of language-specific shorthands (see #6817):
\let\LanguageShortHands\languageshorthands
\def\languageshorthands#1{}
\setlength{\emergencystretch}{3em} % prevent overfull lines
\providecommand{\tightlist}{%
  \setlength{\itemsep}{0pt}\setlength{\parskip}{0pt}}
% Ajouté au préambule de chaque document par generer.sh
\usepackage{xcolor}
\usepackage{newunicodechar}
\newfontfamily\symbolfont{DejaVuSans.ttf}
% Caractères absents de Latin Modern (flèches, comparateurs, cases) : repli sur DejaVu Sans
\newunicodechar{→}{{\symbolfont →}}
\newunicodechar{←}{{\symbolfont ←}}
\newunicodechar{⇄}{{\symbolfont ⇄}}
\newunicodechar{↔}{{\symbolfont ↔}}
\newunicodechar{≥}{{\symbolfont ≥}}
\newunicodechar{≤}{{\symbolfont ≤}}
\newunicodechar{≠}{{\symbolfont ≠}}
\newunicodechar{≈}{{\symbolfont ≈}}
\newunicodechar{✓}{{\symbolfont ✓}}
\newunicodechar{✗}{{\symbolfont ✗}}
\newunicodechar{⚠}{{\symbolfont ⚠}}
\newunicodechar{✅}{{\symbolfont\color[rgb]{0.1,0.55,0.2}✔}}
\newunicodechar{‑}{\mbox{-}}
% Retour à la ligne dans les blocs de code et les URL longues
\usepackage{fvextra}
\DefineVerbatimEnvironment{Highlighting}{Verbatim}{breaklines,commandchars=\\\{\}}
\usepackage{xurl}
\setlength{\emergencystretch}{3em}
\AtBeginEnvironment{longtable}{\small}
% Figures flottantes (pas de demi-pages vides), jamais plus hautes que 85 % de la page
\usepackage{float}
\floatplacement{figure}{htbp}
\usepackage{adjustbox}
\providecommand*\pandocbounded[1]{#1}% absent des documents sans image
\renewcommand*\pandocbounded[1]{\adjustbox{max width=\linewidth,max totalheight=0.85\textheight}{#1}}
\usepackage[font=small,labelfont=bf]{caption}
% Contenu sur une nouvelle page après la table des matières
\let\tdmorigine\tableofcontents
\renewcommand{\tableofcontents}{\tdmorigine\newpage}
\usepackage{fancyhdr}
\pagestyle{fancy}
\fancyhf{}
\fancyhead[L]{\small SmartRecruit — Qualité logicielle ISO/IEC 25010}
\fancyhead[R]{\small\thepage}
\renewcommand{\headrulewidth}{0.4pt}
\usepackage{bookmark}
\IfFileExists{xurl.sty}{\usepackage{xurl}}{} % add URL line breaks if available
\urlstyle{same}
\hypersetup{
  pdftitle={Fiche d\textquotesingle étude --- Membre 4 : sécurité et qualité de code},
  pdflang={fr},
  colorlinks=true,
  linkcolor={Maroon},
  filecolor={Maroon},
  citecolor={Blue},
  urlcolor={Blue},
  pdfcreator={LaTeX via pandoc}}

\title{Fiche d\textquotesingle étude --- Membre 4 : sécurité et qualité
de code}
\author{}
\date{}

\begin{document}
\maketitle

\renewcommand*\contentsname{Table des matières}
{
\hypersetup{linkcolor=}
\setcounter{tocdepth}{3}
\tableofcontents
}
\begin{quote}
Pour : \emph{(à compléter)}. Complète le \href{SOUTENANCE.pdf}{guide de
soutenance} (questions-réponses, membre 4) et la
\href{REPARTITION_TESTS.pdf}{répartition des tests}.
\end{quote}

\section{1. Ta mission en soutenance}\label{1-ta-mission-en-soutenance}

Tu dois pouvoir \textbf{expliquer} les mécanismes de sécurité,
\textbf{relancer} les analyses, \textbf{montrer} les rapports
avant/après et \textbf{répondre} sur :

\begin{itemize}
\tightlist
\item
  \textbf{6.6 Sécurité} : authentification robuste, contrôle
  d\textquotesingle accès par rôle, absence de failles connues (OWASP
  Top 10), avec les couches SAST, SCA et DAST ;
\item
  \textbf{6.7 Maintenabilité} (analyse statique) : complexité, lint,
  structure en couches.
\end{itemize}

\section{2. Ce que tu dois apprendre
(notions)}\label{2-ce-que-tu-dois-apprendre-notions}

\begin{longtable}[]{@{}
  >{\raggedright\arraybackslash}p{(\linewidth - 2\tabcolsep) * \real{0.3333}}
  >{\raggedright\arraybackslash}p{(\linewidth - 2\tabcolsep) * \real{0.6667}}@{}}
\toprule\noalign{}
\begin{minipage}[b]{\linewidth}\raggedright
Notion
\end{minipage} & \begin{minipage}[b]{\linewidth}\raggedright
Ce que tu dois savoir dire
\end{minipage} \\
\midrule\noalign{}
\endhead
\bottomrule\noalign{}
\endlastfoot
SAST & Analyse statique : lit le code source sans
l\textquotesingle exécuter pour trouver des motifs dangereux (Bandit
pour Python, Semgrep multi-langages). \\
SCA & Analyse des dépendances : compare les versions des bibliothèques
aux bases de CVE (pip-audit, npm audit, Trivy sur les images Docker). \\
DAST & Analyse dynamique : attaque l\textquotesingle application en
cours d\textquotesingle exécution, de l\textquotesingle extérieur (OWASP
ZAP). \\
CVE, gravité & Identifiant public d\textquotesingle une vulnérabilité
connue ; gravité LOW → CRITICAL (score CVSS). On a filtré HIGH et
CRITICAL avec Trivy. \\
Faux positif / risque accepté & Alerte non applicable, ou risque connu
et justifié par écrit (nos 3 alertes ZAP). \\
OWASP Top 10 (2021) & Les 10 catégories de failles web les plus
courantes : A01 contrôle d\textquotesingle accès, A02 cryptographie, A03
injection, A05 mauvaise configuration, A06 composants vulnérables, A07
authentification, A09 journalisation\ldots{} \\
JWT & Jeton signé contenant l\textquotesingle identité et le rôle, avec
une date d\textquotesingle expiration ; le serveur vérifie la signature,
il ne fait pas confiance au contenu sans elle. \\
RBAC & Contrôle d\textquotesingle accès par rôle (\texttt{CANDIDATE},
\texttt{RECRUITER}, \texttt{ADMIN}) ; chez nous \texttt{@PreAuthorize}
sur chaque point d\textquotesingle API. \\
HMAC & Signature d\textquotesingle un message avec une clé secrète
partagée : le gateway signe les en-têtes d\textquotesingle identité,
chaque service vérifie ; un appel direct sans signature reçoit 401. \\
BCrypt & Hachage lent et salé des mots de passe : impossible de
retrouver le mot de passe, coûteux à attaquer par force brute. \\
XSS et assainissement & Injection de script dans une page ; on nettoie
toutes les chaînes JSON entrantes (jsoup). \\
Injection SQL & Requête construite avec des données utilisateur ; évitée
par JPA/JPQL paramétré. \\
Rate limiting, verrouillage & Limiter le nombre de requêtes (429) et
verrouiller un compte après plusieurs échecs : contre la force brute. \\
En-têtes de sécurité, CSP & Le navigateur applique des règles envoyées
par le serveur (Content-Security-Policy, etc.) ; ZAP vérifie leur
présence. \\
Complexité cyclomatique (CCN) & Nombre de chemins indépendants dans une
fonction ; au-delà de 15 (notre seuil lizard), elle devient difficile à
tester et à maintenir. \\
Lint & Outil qui signale le code suspect ou non conforme aux règles
(ruff pour Python, oxlint pour TypeScript). \\
\end{longtable}

\section{3. Fichiers à lire, dans
l\textquotesingle ordre}\label{3-fichiers-uxe0-lire-dans-lordre}

\begin{longtable}[]{@{}
  >{\raggedright\arraybackslash}p{(\linewidth - 4\tabcolsep) * \real{0.0115}}
  >{\raggedright\arraybackslash}p{(\linewidth - 4\tabcolsep) * \real{0.7864}}
  >{\raggedright\arraybackslash}p{(\linewidth - 4\tabcolsep) * \real{0.2020}}@{}}
\toprule\noalign{}
\begin{minipage}[b]{\linewidth}\raggedright
\#
\end{minipage} & \begin{minipage}[b]{\linewidth}\raggedright
Fichier
\end{minipage} & \begin{minipage}[b]{\linewidth}\raggedright
Ce que tu y cherches
\end{minipage} \\
\midrule\noalign{}
\endhead
\bottomrule\noalign{}
\endlastfoot
1 &
\texttt{backend/}\allowbreak{}\texttt{api-}\allowbreak{}\texttt{gateway/}\allowbreak{}\texttt{src/}\allowbreak{}\texttt{main/}\allowbreak{}\texttt{java/}\allowbreak{}\texttt{com/}\allowbreak{}\texttt{smartrecruit/}\allowbreak{}\texttt{gateway/}\allowbreak{}\texttt{security/}\allowbreak{}\texttt{Jwt}\allowbreak{}\texttt{Authentication}\allowbreak{}\texttt{Gateway}\allowbreak{}\texttt{Filter.}\allowbreak{}\texttt{java}
& Vérification du JWT à l\textquotesingle entrée (401 si absent ou
invalide). \\
2 &
\texttt{…/}\allowbreak{}\texttt{gateway/}\allowbreak{}\texttt{security/}\allowbreak{}\texttt{Internal}\allowbreak{}\texttt{Identity}\allowbreak{}\texttt{Headers.}\allowbreak{}\texttt{java}
et
\texttt{Rate}\allowbreak{}\texttt{Limiting}\allowbreak{}\texttt{Filter.}\allowbreak{}\texttt{java}
& Signature HMAC des en-têtes d\textquotesingle identité ; limite de 10
requêtes / 5 min / IP (429). \\
3 &
\texttt{backend/}\allowbreak{}\texttt{application-}\allowbreak{}\texttt{service/}\allowbreak{}\texttt{src/}\allowbreak{}\texttt{main/}\allowbreak{}\texttt{java/}\allowbreak{}\texttt{com/}\allowbreak{}\texttt{smartrecruit/}\allowbreak{}\texttt{application/}\allowbreak{}\texttt{security/}\allowbreak{}\texttt{Internal}\allowbreak{}\texttt{Identity}\allowbreak{}\texttt{Verifier.}\allowbreak{}\texttt{java}
et
\texttt{config/}\allowbreak{}\texttt{Security}\allowbreak{}\texttt{Config.}\allowbreak{}\texttt{java}
& Côté service : vérification de la signature, règles par rôle. \\
4 &
\texttt{…/}\allowbreak{}\texttt{application/}\allowbreak{}\texttt{config/}\allowbreak{}\texttt{Jackson}\allowbreak{}\texttt{Sanitization}\allowbreak{}\texttt{Config.}\allowbreak{}\texttt{java}
et
\texttt{security/}\allowbreak{}\texttt{Html}\allowbreak{}\texttt{Sanitizer.}\allowbreak{}\texttt{java}
& Assainissement XSS de toutes les chaînes JSON entrantes. \\
5 &
\texttt{backend/}\allowbreak{}\texttt{auth-}\allowbreak{}\texttt{service/}\allowbreak{}\texttt{src/}\allowbreak{}\texttt{main/}\allowbreak{}\texttt{java/}\allowbreak{}\texttt{com/}\allowbreak{}\texttt{smartrecruit/}\allowbreak{}\texttt{auth/}\allowbreak{}\texttt{service/}\allowbreak{}\texttt{Auth}\allowbreak{}\texttt{Service.}\allowbreak{}\texttt{java}
et
\texttt{config/}\allowbreak{}\texttt{Security}\allowbreak{}\texttt{Config.}\allowbreak{}\texttt{java}
& Hachage BCrypt, verrouillage après 5 échecs en 15 min (RG‑06). \\
6 &
\texttt{backend/}\allowbreak{}\texttt{user-}\allowbreak{}\texttt{service/}\allowbreak{}\texttt{src/}\allowbreak{}\texttt{main/}\allowbreak{}\texttt{java/}\allowbreak{}\texttt{com/}\allowbreak{}\texttt{smartrecruit/}\allowbreak{}\texttt{user/}\allowbreak{}\texttt{storage/}\allowbreak{}\texttt{Cv}\allowbreak{}\texttt{Cipher.}\allowbreak{}\texttt{java}
& Chiffrement des CV stockés. \\
7 &
\texttt{backend/}\allowbreak{}\texttt{*/}\allowbreak{}\texttt{src/}\allowbreak{}\texttt{test/}\allowbreak{}\texttt{java/}\allowbreak{}\texttt{…/}\allowbreak{}\texttt{security/}\allowbreak{}\texttt{Rbac}\allowbreak{}\texttt{Integration}\allowbreak{}\texttt{Test.}\allowbreak{}\texttt{java}
et
\texttt{backend/}\allowbreak{}\texttt{api-}\allowbreak{}\texttt{gateway/}\allowbreak{}\texttt{src/}\allowbreak{}\texttt{test/}\allowbreak{}\texttt{java/}\allowbreak{}\texttt{…/}\allowbreak{}\texttt{Rate}\allowbreak{}\texttt{Limiting}\allowbreak{}\texttt{Integration}\allowbreak{}\texttt{Test.}\allowbreak{}\texttt{java}
& Les tests 401/403 et 429. \\
8 & \texttt{quality/run-all.sh}, étapes \texttt{sast}, \texttt{sca},
\texttt{dast}, \texttt{maintainability} & Quels outils, quelles règles
(Semgrep : \texttt{p/owasp-top-ten}, \texttt{p/java},
\texttt{p/typescript}, \texttt{p/python}, \texttt{p/secrets}), quels
fichiers exclus. \\
9 &
\texttt{quality/}\allowbreak{}\texttt{reports/}\allowbreak{}\texttt{security/}\allowbreak{}
: \texttt{trivy-resume.txt}, \texttt{*-avant*}, \texttt{zap/} & Comparer
avant/après correction. \\
\end{longtable}

\section{4. Commandes à savoir lancer (et
lire)}\label{4-commandes-uxe0-savoir-lancer-et-lire}

\begin{Shaded}
\begin{Highlighting}[]
\FunctionTok{bash}\NormalTok{ quality/run{-}all.sh sast}
\end{Highlighting}
\end{Shaded}

Lire
\texttt{quality/}\allowbreak{}\texttt{reports/}\allowbreak{}\texttt{security/}\allowbreak{}\texttt{bandit.}\allowbreak{}\texttt{txt}
(aucun problème) et le résumé de \texttt{semgrep.txt}.

\begin{Shaded}
\begin{Highlighting}[]
\FunctionTok{bash}\NormalTok{ quality/run{-}all.sh sca}
\end{Highlighting}
\end{Shaded}

Images Docker construites au préalable. Lire \texttt{pip-audit.txt},
\texttt{npm-audit.txt}, puis \texttt{trivy-resume.txt} (0
HIGH/CRITICAL).

\begin{Shaded}
\begin{Highlighting}[]
\FunctionTok{bash}\NormalTok{ quality/run{-}all.sh dast}
\end{Highlighting}
\end{Shaded}

Application démarrée. Ouvrir
\texttt{quality/}\allowbreak{}\texttt{reports/}\allowbreak{}\texttt{security/}\allowbreak{}\texttt{zap/}\allowbreak{}\texttt{zap-}\allowbreak{}\texttt{frontend.}\allowbreak{}\texttt{html}
: FAIL, WARN, PASS.

\begin{Shaded}
\begin{Highlighting}[]
\FunctionTok{bash}\NormalTok{ quality/run{-}all.sh maintainability}
\end{Highlighting}
\end{Shaded}

Lire les dernières lignes de
\texttt{quality/}\allowbreak{}\texttt{reports/}\allowbreak{}\texttt{maintainability/}\allowbreak{}\texttt{lizard-}\allowbreak{}\texttt{complexite.}\allowbreak{}\texttt{txt}
(CCN moyen, fonctions au-dessus du seuil), puis \texttt{ruff.txt} et
\texttt{oxlint.txt}.

\section{5. Chiffres à connaître par
cœur}\label{5-chiffres-uxe0-connauxeetre-par-cux153ur}

\begin{longtable}[]{@{}
  >{\raggedright\arraybackslash}p{(\linewidth - 2\tabcolsep) * \real{0.2632}}
  >{\raggedright\arraybackslash}p{(\linewidth - 2\tabcolsep) * \real{0.7368}}@{}}
\toprule\noalign{}
\begin{minipage}[b]{\linewidth}\raggedright
Sujet
\end{minipage} & \begin{minipage}[b]{\linewidth}\raggedright
Valeur
\end{minipage} \\
\midrule\noalign{}
\endhead
\bottomrule\noalign{}
\endlastfoot
SAST & \textbf{0} problème ; Semgrep : \textbf{355} règles sur
\textbf{498} fichiers \\
SCA images (Trivy) & \textbf{179 → 0} CVE HIGH/CRITICAL, dont
\textbf{24} critiques, sur \textbf{8} images \\
SCA Python / npm & \textbf{82 → 0} (pip-audit), \textbf{1 → 0} (npm
audit) \\
DAST (ZAP) & \textbf{0} échec, \textbf{59} contrôles réussis, \textbf{3}
alertes moyennes acceptées \\
Anti-force brute & verrouillage après \textbf{5} échecs en \textbf{15
min} ; \textbf{10} requêtes / \textbf{5 min} / IP sur connexion et
inscription, puis \textbf{429} \\
Complexité & CCN moyen \textbf{2,2} ; \textbf{5} fonctions sur
\textbf{875} au-dessus de 15 \\
Lint & ruff et oxlint : \textbf{0} avertissement, bloquants en CI \\
\end{longtable}

\section{6. Preuves à savoir ouvrir}\label{6-preuves-uxe0-savoir-ouvrir}

\begin{itemize}
\tightlist
\item
  \texttt{quality/}\allowbreak{}\texttt{reports/}\allowbreak{}\texttt{security/}\allowbreak{}\texttt{bandit.}\allowbreak{}\texttt{txt},
  \texttt{semgrep.txt}.
\item
  \texttt{quality/}\allowbreak{}\texttt{reports/}\allowbreak{}\texttt{security/}\allowbreak{}\texttt{pip-}\allowbreak{}\texttt{audit-}\allowbreak{}\texttt{avant.}\allowbreak{}\texttt{txt}
  puis \texttt{pip-audit.txt} ; \texttt{trivy-avant/} puis
  \texttt{trivy/} et \texttt{trivy-resume.txt}.
\item
  \texttt{quality/}\allowbreak{}\texttt{reports/}\allowbreak{}\texttt{security/}\allowbreak{}\texttt{zap/}\allowbreak{}\texttt{zap-}\allowbreak{}\texttt{frontend.}\allowbreak{}\texttt{html}
  (et \texttt{zap-avant/} pour comparer).
\item
  \texttt{quality/}\allowbreak{}\texttt{reports/}\allowbreak{}\texttt{security/}\allowbreak{}\texttt{rate-}\allowbreak{}\texttt{limit-}\allowbreak{}\texttt{429.}\allowbreak{}\texttt{txt}
  : preuve d\textquotesingle exécution réelle de la limite.
\item
  \texttt{quality/}\allowbreak{}\texttt{reports/}\allowbreak{}\texttt{maintainability/}\allowbreak{}\texttt{lizard-}\allowbreak{}\texttt{complexite.}\allowbreak{}\texttt{txt},
  \texttt{ruff.txt}, \texttt{oxlint.txt}.
\end{itemize}

\section{7. Questions pièges (au-delà du guide de
soutenance)}\label{7-questions-piuxe8ges-au-deluxe0-du-guide-de-soutenance}

\begin{longtable}[]{@{}
  >{\raggedright\arraybackslash}p{(\linewidth - 2\tabcolsep) * \real{0.4286}}
  >{\raggedright\arraybackslash}p{(\linewidth - 2\tabcolsep) * \real{0.5714}}@{}}
\toprule\noalign{}
\begin{minipage}[b]{\linewidth}\raggedright
Question
\end{minipage} & \begin{minipage}[b]{\linewidth}\raggedright
Réponse
\end{minipage} \\
\midrule\noalign{}
\endhead
\bottomrule\noalign{}
\endlastfoot
Où est stocké le JWT et que se passe-t-il s\textquotesingle il est volé
? & Dans le \texttt{localStorage} du navigateur
(\texttt{frontend/}\allowbreak{}\texttt{src/}\allowbreak{}\texttt{api/}\allowbreak{}\texttt{client.}\allowbreak{}\texttt{ts}).
Volé, il reste valable jusqu\textquotesingle à son expiration :
\textbf{1 h} pour le jeton d\textquotesingle accès (7 jours pour le
jeton de renouvellement,
\texttt{auth-}\allowbreak{}\texttt{service/}\allowbreak{}\texttt{…/}\allowbreak{}\texttt{application.}\allowbreak{}\texttt{yml}).
Le principal risque de vol est le XSS, d\textquotesingle où
l\textquotesingle assainissement des entrées et la CSP. Un cookie
\texttt{HttpOnly} serait une amélioration possible. \\
Pourquoi des en-têtes signés en plus du JWT ? & Le JWT est vérifié une
fois au gateway. Sans signature, n\textquotesingle importe qui accédant
au réseau interne pourrait appeler un service en se prétendant admin. La
signature HMAC garantit que l\textquotesingle identité vient du gateway,
et une signature de plus de \textbf{2 minutes} est refusée
(\texttt{Internal}\allowbreak{}\texttt{Identity}\allowbreak{}\texttt{Verifier}),
ce qui empêche de rejouer une requête capturée. \\
SonarQube était proposé : pourquoi ne pas l\textquotesingle utiliser ? &
Aucun outil n\textquotesingle est imposé (§7). Semgrep (sécurité),
lizard (complexité), ruff/oxlint (lint) et JaCoCo/pytest-cov/Vitest
(couverture) couvrent les mêmes rôles, s\textquotesingle exécutent en
local et en CI sans serveur à héberger. \\
Comment avez-vous corrigé 179 CVE ? & Montées de version des dépendances
et des images de base (Spring Boot 4.1.2, jsoup 1.24.1,
FastAPI/starlette, pypdf\ldots), pip retiré de l\textquotesingle image
IA, puis tous les tests rejoués pour vérifier qu\textquotesingle aucune
régression n\textquotesingle apparaissait. \\
0 vulnérabilité, cela veut dire que l\textquotesingle application est
sûre ? & Non : 0 vulnérabilité \textbf{connue détectée} par ces outils.
Ils ne trouvent pas les failles de logique métier ;
c\textquotesingle est pourquoi les règles d\textquotesingle accès sont
aussi testées (401/403, RBAC, \texttt{/forbidden}). \\
Pourquoi accepter 3 alertes ZAP ? & Elles sont justifiées :
\texttt{unsafe-inline} nécessaire pour les styles React, images
d\textquotesingle entreprise hébergées à l\textquotesingle extérieur,
pas de SRI possible sur Google Fonts (feuille générée dynamiquement).
Accepter un risque documenté est une décision, pas un oubli. \\
Les 5 fonctions complexes, un problème ? & 5 sur 875, au-dessus du seuil
de 15 : signalées et toutes testées ; candidates à un découpage si elles
évoluent. \\
La sécurité des CV ? & Fichiers chiffrés au stockage
(\texttt{CvCipher}), consultation journalisée dans
l\textquotesingle audit (\texttt{CV\_VIEWED}), effacement RGPD
possible. \\
\end{longtable}

\section{8. Plan
d\textquotesingle apprentissage}\label{8-plan-dapprentissage}

\begin{longtable}[]{@{}
  >{\raggedright\arraybackslash}p{(\linewidth - 4\tabcolsep) * \real{0.1316}}
  >{\raggedright\arraybackslash}p{(\linewidth - 4\tabcolsep) * \real{0.1316}}
  >{\raggedright\arraybackslash}p{(\linewidth - 4\tabcolsep) * \real{0.7368}}@{}}
\toprule\noalign{}
\begin{minipage}[b]{\linewidth}\raggedright
Étape
\end{minipage} & \begin{minipage}[b]{\linewidth}\raggedright
Durée
\end{minipage} & \begin{minipage}[b]{\linewidth}\raggedright
Objectif
\end{minipage} \\
\midrule\noalign{}
\endhead
\bottomrule\noalign{}
\endlastfoot
1 & 2 h & Apprendre les notions de la section 2, en priorité
l\textquotesingle OWASP Top 10 2021 (page officielle) et le
fonctionnement d\textquotesingle un JWT. \\
2 & 2 h & Lire les fichiers 1 à 7 de la section 3 ; suivre une requête
du navigateur au service (figure 7 du rapport de conception). \\
3 & 1 h & Lancer \texttt{sast}, \texttt{sca}, \texttt{maintainability} ;
ouvrir chaque preuve de la section 6. \\
4 & 1 h & Lancer \texttt{dast} (application démarrée) et parcourir le
rapport ZAP ; comparer avec \texttt{zap-avant/}. \\
5 & 1 h & Réciter les réponses du guide de soutenance (membre 4) et de
la section 7. \\
\end{longtable}

\section{9. Auto-évaluation}\label{9-auto-uxe9valuation}

\begin{itemize}
\tightlist
\item[$\square$]
  J\textquotesingle explique SAST, SCA et DAST avec
  l\textquotesingle outil que nous utilisons pour chacun.
\item[$\square$]
  Je suis le chemin d\textquotesingle une requête : JWT au gateway →
  en-têtes HMAC → vérification dans le service → \texttt{@PreAuthorize}.
\item[$\square$]
  Je relie au moins 6 catégories OWASP Top 10 à un mécanisme et à un
  test du projet.
\item[$\square$]
  Je montre un rapport avant/après (Trivy ou pip-audit) et
  j\textquotesingle explique la correction.
\item[$\square$]
  Je justifie les 3 alertes ZAP acceptées.
\item[$\square$]
  J\textquotesingle explique la complexité cyclomatique et je lis la fin
  du rapport lizard.
\end{itemize}

\end{document}
